
\subsubsection{Key Findings}

\textbf{Finding 1: CV\_PR positively correlates with accuracy (r = +0.61).}

The original mechanism---flat spectra lead to stable SVD estimates, indicating smooth loss landscapes enabling generalization---is contradicted. Models with higher accuracy exhibit more variance in participation ratio across random projections.

\textbf{Finding 2: The extraction methodology is reliable.}

100\% completion across 100 diverse models, with CV\_PR values consistently finite and in a narrow range [0.0014, 0.0326], validates that randomized SVD with 20 seeds produces meaningful variance estimates.

\textbf{Finding 3: The falsification redirects research.}

The decisive falsification (r = +0.61 vs predicted r < -0.3) indicates that future work should investigate why higher variance accompanies better models rather than pursuing stability-as-quality mechanisms.

\subsubsection{Alternative Explanations}

\textbf{Confounding by model size.} Larger models tend to have higher accuracy and potentially more spectral diversity. Without partial correlation controlling for parameter count, it is not possible to determine whether CV\_PR has independent predictive value or proxies model size.

\textbf{Richer representations.} Higher CV\_PR may reflect diverse feature extraction at different scales. Better models may learn more varied spectral structures across layers, increasing projection-dependent variance.

\subsubsection{Limitations}

\textbf{Limitation 1: Confounding not controlled.}

This was designed as an existence test. The purpose was to establish whether correlation exists and in which direction, not to fully characterize causal mechanisms. Partial correlation controlling for parameter count would be required to distinguish whether CV\_PR has independent predictive value.

\textbf{Limitation 2: Mechanism hypotheses blocked.}

The MUST\_WORK failure on H-E2 blocked dependent hypotheses (H-M1: partial correlation controlling condition number; H-M2: within-family correlation). The mechanism for why CV\_PR relates to accuracy remains unknown.

\textbf{Limitation 3: Limited architecture stratification.}

While multiple architecture families (ResNet, ViT, EfficientNet, ConvNeXt, etc.) were included, within-family versus cross-family correlation patterns were not analyzed. The relationship may differ by architecture type.

\subsubsection{Broader Impact}

This work contributes to rigorous hypothesis testing in machine learning research. The falsification provides a concrete starting point for reformulated hypotheses. Misinterpretation of results should be avoided---the observed positive correlation may be spurious (driven by confounding) and should not be used for model selection without controlling for model size.
